%% kapa: exam=1 round=33 year=2022 period=1 form=A subject=economics subject_ko=경제학원론 seq=19 number=59
%% source: Q-Net 감정평가사 기출문제 (cst00309) artlSeq=5209997 file=raw/1008/5209997/2022년도 제33회 감정평가사 1차시험 1교시(민법,경제학,부동산학).pdf
%% conversion: kordoc 4.24.51 -> tools/gen1.py (auto); review=proofread-sonnet; flags=none
\begin{question}{59}
노동시장이 수요독점일 때 이에 관한 설명으로 옳은 것을 모두 고른 것은? (단, 생산물 시장은 완전경쟁시장이며, 노동수요곡선은 우하향, 노동공급곡선은 우상향한다.)
\begin{framedtext}
ㄱ. 노동의 한계생산가치(value of marginal product of labor)곡선이 노동수요곡선이다. \par
ㄴ. 한계요소비용(marginal factor cost)곡선은 노동공급곡선의 아래쪽에 위치한다. \par
ㄷ. 균형 고용량은 노동의 한계생산가치곡선과 한계요소비용곡선이 만나는 점에서 결정된다. \par
ㄹ. 노동시장이 완전경쟁인 경우보다 균형 임금률이 낮고 균형 고용량이 많다.
\end{framedtext}
\begin{choices}
  \item ㄱ, ㄴ
  \item ㄱ, ㄷ
  \item ㄱ, ㄹ
  \item ㄴ, ㄷ
  \item ㄷ, ㄹ
\end{choices}
\end{question}
\answer{2}
